% uid: ""
% school: "lycee"
% level: "3eme"
% section: ["Economie et Gestion"]
% chapter: "Statistiques"
% title: "Ventes mensuelles et droite de Mayer"
% difficulty: "Moyen"
% points: null
% tags: ["statistiques", "serie-double", "nuage-de-points", "point-moyen", "droite-de-mayer", "prevision"]
% provenance: "own"
% dressing: "droite de Mayer"
% rung: 3
% batch: "2026-11-20-stat-3eco-A"

% ---STATEMENT---
Une boutique en ligne a relevé ses ventes $y$ (en centaines d'articles) pendant les $8$ premiers mois de l'année ($x$ = numéro du mois) :

\begin{center}
\begin{tabular}{|l|c|c|c|c|c|c|c|c|}
\hline
Mois $x_i$ & $1$ & $2$ & $3$ & $4$ & $5$ & $6$ & $7$ & $8$ \\ \hline
Ventes $y_i$ & $14$ & $16$ & $18$ & $20$ & $24$ & $26$ & $28$ & $30$ \\ \hline
\end{tabular}
\end{center}

\begin{enumerate}
  \item Représenter le nuage de points $M_i(x_i\,;y_i)$ et calculer les coordonnées du point moyen $G$.
  \item On partage le nuage en deux sous-nuages : les $4$ premiers mois et les $4$ derniers. Calculer les points moyens $G_1$ et $G_2$.
  \item Déterminer l'équation de la droite de Mayer $(G_1G_2)$ et vérifier qu'elle passe par $G$.
  \item Prévoir les ventes du mois $10$.
\end{enumerate}

% ---SOLUTION---
\begin{enumerate}
  \item Le nuage est formé des points $(1;14)$, $(2;16)$, $\dots$, $(8;30)$ : il est croissant et presque aligné.

  Point moyen : $\overline{x}=\dfrac{1+2+\dots+8}8=\dfrac{36}8=4{,}5$ et $\overline{y}=\dfrac{14+16+18+20+24+26+28+30}8=\dfrac{176}8=22$.
  $$ \ib{$G(4{,}5\,;\,22)$}.$$
  \item \textbf{Rappel.}\par
  \begin{prop}[Méthode de Mayer]
  On range les points par abscisses croissantes, on les partage en deux sous-nuages de même effectif, on calcule leurs points moyens $G_1$ et $G_2$ : la droite d'ajustement est $(G_1G_2)$.
  \end{prop}
  Mois $1$ à $4$ : $\overline{x}_1=\dfrac{1+2+3+4}4=2{,}5$, $\overline{y}_1=\dfrac{14+16+18+20}4=\dfrac{68}4=17$ : $G_1(2{,}5\,;17)$.

  Mois $5$ à $8$ : $\overline{x}_2=\dfrac{5+6+7+8}4=6{,}5$, $\overline{y}_2=\dfrac{24+26+28+30}4=\dfrac{108}4=27$ : $G_2(6{,}5\,;27)$.
  \item Pente : $a=\dfrac{y_{G_2}-y_{G_1}}{x_{G_2}-x_{G_1}}=\dfrac{27-17}{6{,}5-2{,}5}=\dfrac{10}4=2{,}5$.

  La droite passe par $G_1$ : $17=2{,}5\times2{,}5+b$, donc $b=17-6{,}25=10{,}75$.
  $$ \ib{$y=2{,}5\,x+10{,}75$}.$$
  \emph{Vérification} : $2{,}5\times4{,}5+10{,}75=11{,}25+10{,}75=22=\overline{y}$ $\checkmark$. $G$ appartient à la droite, ce qui est normal : $G$ est le milieu de $[G_1G_2]$ (même effectif dans les deux sous-nuages).
  \item Pour $x=10$ : $y=2{,}5\times10+10{,}75=35{,}75$ : environ $3\,575$ articles.

  \textbf{Piège :} $x=10$ est proche des données ($1\le x\le8$), la prévision est raisonnable ; on se méfierait d'une prévision à $x=40$.
\end{enumerate}

\textbf{Idée :} la droite de Mayer passe par les deux points moyens des sous-nuages, donc aussi par le point moyen du nuage.
